% Part: conditional-logics
% Chapter: minimal-change-semantics
% Section: sphere-models

\documentclass[../../../include/open-logic-section]{subfiles}

\begin{document}

\olfileid{con}{min}{sph}

\olsection{د کروي سيمو مدلونه}

د واقعيت خلاف شرطونو د رسمي معنٰی پوهنې د ورکولو يوه طريقه دا
ده چې د لويس غيررسمي شننه په رياضيکي جوړښت بدله کړو۔ بيا د
يوې نړۍ~$w$ چاپېره کروي سيمې د نړۍو سټونه دي۔ څرنګه چې
کروي سيمې يو په بل کښې دننه دي، د~$w$ چاپېره د نړۍو سټونه
بايد د فرعي سټ د اړيکې له مخې خطي ترتيب ولري۔

\begin{defn}
  \emph{د کروي سيمو مدل} يو درې‌توب~$\mModel{M} = \tuple{W, O, V}$
  دے، چې $W$ د نړۍو يو ناتش سټ، $V\colon \PVar \to \Pow{W}$
  يوه ارزښت ټاکنه، او $O\colon W \to \Pow{\Pow{W}}$ داسې تابع ده
  چې هرې نړۍ~$w$ ته \emph{د کروي سيمو يو نظام}~$O_w$ ټاکي۔
  د هر $w$ لپاره، $O_w$ د نړۍو د سټونو يو سټ دے او بايد دا شرطونه پوره کړي:
  \begin{enumerate}
  \item $O_w$ پر~$w$ \emph{مرکوز} دے: $\{w\} \in O_w$۔
  \item $O_w$ \emph{يو په بل کښې دننه} دے: هر کله چې $S_1$، $S_2 \in O_w$، نو $S_1
    \subseteq S_2$ يا $S_2 \subseteq S_1$؛ يعنې $O_w$ د~$\subseteq$
    له مخې خطي ترتيب لري۔
  \item $O_w$ د ناتشو کورنيو د اتحاد پر وړاندې تړلے دے۔
  \item $O_w$ د ناتشو کورنيو د تقاطع پر وړاندې تړلے دے۔
  \end{enumerate}
\end{defn}

د $O_w$ تر شا تصور دا دے چې د $w$ «چاپېره» نړۍ له~$w$ څخه
د واټن له مخې په پوړونو کښې وېشل شوې دي۔ تر ټولو دننۍ ناتشه
کروي سيمه يوازې $w$، يعنې سټ~$\{w\}$، ده: $w$ له~$w$ سره
تر هغو نړۍو نژدې دے چې په بله کروي سيمه کښې دي۔ د نسخې
څرګندونه (OLCNT-008): منجمد رسمي تعريف تش سټ نۀ منع کوي؛
دلته د تر ټولو دننۍ سيمې بيان د ناتشو سيمو په معنا روښانه شوے
دے، رسمي تعريف عين دے۔ که $S \subsetneq S'$، نو په $S' \setminus S$
کښې نړۍ له $w$ څخه تر هغو نړۍو ليرې دي چې په~$S$ کښې دي:
$S' \setminus S$ د $S$ او د~$S'$ د بهرنۍ سطحې ترمنځ «پوړ»
دے۔ د نسخې څرګندونه (OLCNT-006): منجمده سرچينه دا د دننۍ
سيمې او د بهرنۍ سيمې څخه بهر نړۍو ترمنځ پوړ بولي؛ د ليکل شوي
سټ توپير له مخې دغه پوړ له دننۍ سيمې بهر خو په بهرنۍ سيمه
کښې دننه دے۔ په تېره بيا، بايد کروي سيمې داسې وګڼو چې د خپلې
بهرنۍ سطحې دننه ټولې نړۍ رانغاړي؛ يوازې جلا پوړونه نۀ دي۔

\begin{figure}
\begin{center}
\begin{tikzpicture}[layerwidth=1.5,scale=.6]\tiny
  \spheresystem{3}
  \propositionintersect{0}{3}{60}{6}
  \path (0,0) node[world] {$w$};
  \spherepos{-30}{2}{node[world] {$w_2$}}
  \spherepos{220}{2}{node[world] {$w_3$}}
  \spherepos{90}{2}{node[world] {$w_1$}}
  \spherepos[shift={(.1,0)}]{10}{3}{node[world] {$w_5$}}
  \spherepos[shift={(.1,0)}]{-10}{3}{node[world] {$w_6$}}
  \spherepos{225}{3}{node[world] {$w_4$}}
  \spherepos{20}{4}{node[world] {$w_7$}}
  \path (5,0) node[left] {$p$};
\end{tikzpicture}
\caption{د کروي سيمو د يو مدل انځور}
\ollabel{fig:sphere-model}
\end{center}
\end{figure}

په \olref{fig:sphere-model} کښې انځور له هغه د کروي سيمو له
مدل سره برابر دے چې $W = \{w, w_1, \dots, w_7\}$ او
$V(p) = \{w_5, w_6,
w_7\}$ لري۔ تر ټولو دننۍ کروي سيمه $S_1 = \{w\}$ ده۔
له $w$ سره تر ټولو نژدې نړۍ $w_1, w_2, w_3$ دي، نو ورپسې
لويه کروي سيمه $S_2 = \{w, w_1,
w_2, w_3\}$ ده۔ نورې ليرې نړۍ $w_4$، $w_5$ او $w_6$ دي،
نو تر ټولو بهرنۍ کروي سيمه $S_3 = \{w, w_1, \dots, w_6\}$
ده۔ د $w$ چاپېره د کروي سيمو نظام $O_w = \{S_1, S_2, S_3\}$
دے۔ نړۍ~$w_7$ د~$w$ چاپېره په هېڅ کروي سيمه کښې نۀ ده۔
هغه تر ټولو نژدې نړۍ چې $p$~پکښې رښتينے دے $w_5$ او~$w_6$
دي، نو تر ټولو کوچنۍ $p$-منونکې کروي سيمه~$S_3$ ده۔

د کروي سيمو په مدل~$\mModel M$ کښې په نړۍ~$w$ کښې د
فورمول~$!A$ تحقق، $\mSat{M}{!A}[w]$، د تعريفولو لپاره د
مودالي !!{formula}s تعريف داسې غځوو چې د $!B \cif !C$ لپاره
يو شرط ورشامل شي:

\begin{defn}
  $\mSat{M}{!B \cif !C}[w]$ هله او يوازې هله چې له دې دواړو يو وي:
  \begin{enumerate}
  \item\ollabel{sphere-vac} د ټولو $u \in \bigcup O_w$ لپاره، $\mSat/{M}{!B}[u]$، يا
  \item\ollabel{sphere-nonvac} د يوې $S \in O_w$ لپاره،
    \begin{enumerate}
      \item د يوه $u \in
        S$ لپاره $\mSat{M}{!B}[u]$، او
      \item د ټولو $v \in S$ لپاره، يا
        $\mSat/{M}{!B}[v]$ يا $\mSat{M}{!C}[v]$۔
    \end{enumerate}
  \end{enumerate}
\end{defn}

د دې تعريف له مخې، $\mSat{M}{!B \cif !C}[w]$ هله او يوازې
هله چې يا مقدم~$!B$ د~$w$ چاپېره په کروي سيمو کښې هر چېرته
کاذب وي، يا يوه کروي سيمه~$S$ وي چې $!B$ پکښې رښتينے وي او
مادي شرطي $!B \lif !C$ په دغه «$!B$-منونکې» کروي سيمه کښې
په ټولو نړۍو کښې رښتينې وي۔ پام وکړئ چې په تعريف کښې مو دا
شرط نۀ دے ايښے چې $S$ دې \emph{تر ټولو دننۍ} $!B$-منونکې
کروي سيمه وي، که څه هم د لومړني تصور له تشريح داسې تمه کېدے
شي۔ خو که په~\olref{sphere-nonvac} کښې شرط د يوې کروي
سيمې~$S$ لپاره پوره وي، نو د هرې هغې کروي سيمې لپاره هم پوره
دے چې~$S$ يې رانغاړي او مقدم پکښې په کومه نړۍ کښې رښتينے
وي؛ نو په تېره بيا د تر ټولو دننۍ منونکې سيمې لپاره، که
موجوده وي۔ د نسخې څرګندونه (OLCNT-007): منجمده سرچينه د ټولو
دننيو سيمو دعوه کوي، خو د شاهد د شتون لومړے فرعي شرط په
نۀ منونکې دننۍ سيمه کښې نۀ ساتل کېږي؛ دلته د منونکې سيمې قيد
څرګند شوے دے۔

دا هم په پام کښې ونيسئ چې د کروي سيمو د مدلونو تعريف دا
شرط نۀ ږدي چې يوه تر ټولو دننۍ $!B$-منونکې کروي سيمه
\emph{موجوده وي}: کېدے شي د $!B$-منونکو کروي سيمو يو
نامتناهي لړ $S_1 \supsetneq S_2 \supsetneq \dots \supsetneq
\{w\}$ ولرو او په دې توګه تر ټولو دننۍ $!B$-منونکې کروي
سيمه ونۀ لرو۔ په دې حالت کښې، $\mSat{M}{!B \cif !C}[w]$
هله او يوازې هله چې $!B \lif !C$ د يوه~$i$ لپاره په ټولو
کروي سيمو $S_i$، $S_{i+1}$، \dots کښې برقرار وي۔

\end{document}
